\section{Discussion}

\subsection{Interpretation}

The central hypothesis concerns the evidence set, not the complexity of the graph encoder. A positive result requires an advantage under valid future-event evidence or a measurable reduction in sensitivity to copied sources. If parameter-matched cluster pooling performs as well as the graph model, independent-source aggregation remains relevant but the graph contribution must be removed. If cutoff filtering helps only when the unrestricted baseline receives label-revealing future documents, the study supports a leakage warning rather than a general accuracy claim.

\subsection{Ethical Considerations}

The corpus inherits professional verdicts and may reproduce institutional or topical biases in the selected fact-checkers. The system is intended to assist review, not to assign reputational penalties or automate content removal. Public release is limited to code, URLs, hashes, timestamps, derived annotations, and split manifests unless source licenses permit redistribution. Search queries and archives may expose deleted or sensitive material; collection logs therefore require access controls and a documented takedown procedure.

\subsection{Limitations}

The planned corpus is English-language and United States-specific. A single-author 200-case audit checks extraction and protocol adherence but does not establish inter-annotator reliability. Archive coverage is incomplete, and first observed time is an upper bound on public availability rather than proof of original publication. Dependency clustering cannot always recover the true reporting chain, especially when copying is paraphrased or offline. The 2026 test is a partial-year snapshot. Live web retrieval is nonstationary, and model comparisons may differ in pretraining and tool access even under a shared candidate pool. These limits constrain geographic, temporal, and causal generalization.
